\newpage

\section{\huge \textcolor{blue}{\hspace*{0.2em}线性齐次微分方程组}}
\subsection{\LARGE \hspace*{0.1em}一阶常微分方程组及几何解释}
\subsubsection{\large 线性方程组}
形如:
\[
\begin{cases}
 x' = f(x,y,t) \\
y' = g(x,y,t)
\end{cases}
\]
初始条件: $x(t_0)=x_0,y(t_0)=y_0$\\
其中$x, y$是因变量, $t$是自变量, 称为一阶方程组.\\
线性方程组:
\[
\begin{cases}
x' = ax + by + r_1(t) \\
y' = cx + dy + r_2(t)
\end{cases}
\]
当$a,b,c,d$是常数时, 称为常系数线性方程组\\
若$a,b,c,d$是$t$的函数, 依旧保持线性\\
而线性齐次将作为首先研究对象,即:
\[ r_1 = 0 \qquad r_2 = 0 \]
依旧以温度模型(双层)为例
列出下方程组(弱化物理属性):
\[
\begin{cases}
    \frac{dT_1}{dt} = a (T_2 - T_1)\\
    \frac{dT_2}{dt} = a (T_1 - T_2) + b (T_e - T_2)
\end{cases}
\] 
整理得:
\[
\begin{cases}
    T_1' = -a T_1 + a T_2 \\
    T_2' = a T_1 - (a+b) T_2 + b T_e 
\end{cases}
\]
假设 
$\begin{cases}
    a=2\\
    b=3
\end{cases}$, $T_e=0$, 则:
\[
\begin{cases}
    T_1' = -2 T_1 + 2 T_2 \\
    T_2' = \hspace*{0.6em}2 T_1 - 5 T_2 
\end{cases}
\]
法一:(消去变量)
\begin{align*}
    T_2 &= \frac{T_1'+2T_1}{2} \\
    (\frac{T_1'+2T_1}{2})' &= 2T_1 - 5(\frac{T_1'+2T_1}{2}) \\
    \intertext{整理得:}
    T_1'' + 7 T_1' + 6 T_1 &= 0 \\
    r^2 + 7 r + 6 &= 0 \\
    (r+1)(r+6) &= 0 \\
    r &= -1 \hspace*{0.5em} or -6
    \intertext{解得:}
    T_1 &= c_1 e^{-t} + c_2 e^{-6t}
\end{align*}
为了确保最后的解只有2个常数, 需用已知的$T_1$找$T_2$
\[ T_2 = \frac{T_1'+2T_1}{2} \]
\[ \therefore  \hspace*{1em} T_2 = \frac{1}{2}c_1 e^{-t} - 2c_2 e^{-6t} \]
假设初值
$\begin{cases}
    T_1(0) = 40 \\
    T_2(0) = 45
\end{cases}$, 代入:
\[
\begin{cases}
    40 = c_1 +c_2 \\
    45 = \frac{1}{2}c_1 - 2c_2
\end{cases}
\Longrightarrow  \hspace*{1em}
\begin{cases}
    c_1 = 50 \\
    c_2 = -10
\end{cases}
\]

\subsubsection{\large 几何解释}
求解过后, 试分析其几何意义\\
若方程如下:
\[
\begin{cases}
    x' = f(x,y) \\
    y' = g(x,y)
\end{cases}
\]
在右侧没有$t$的情况下, 称自治方程组\\
解:
\[ 
\begin{cases}
    x = x(t) \\
    y = y(t)
\end{cases}
\]
解为参数方程组的形式, 即形成参数化曲线\\
\begin{figure}[htbp]
    \centering
\includegraphics[width=0.5\textwidth]{C:/Users/37/Desktop/Notes on Common Calculus Equations/Notes and pictures/cwf11.2.1.png}
    \caption{参数化曲线}
    \label{fig:cwf11.2.1}
\end{figure}

图中存在多条曲线, 因为方程存在多解, 当
输入的初始条件不同, 则会生成不同曲线\\
表示:
$\begin{cases}
    x'(t) \\
    y'(t)
\end{cases}$, 为时刻$t$解的速度, 即\\
此自治微分方程组的几何意义为速度场, 解为参数化曲线.

\subsection{\LARGE \hspace*{0.1em}常系数齐次方程组}
\subsubsection{\large 尝试解-法-矩阵形式}
重新审视前一节的温度模型方程组\\
为了让变量更中性, 令$x=T_1, y=T_2$, 则:
\[
\begin{cases}
    x' = -2 x + 2 y \\
    y' = 2 x - 5 y
\end{cases}
\]
法一: 消元, 前文已介绍, 解为:
\[
\begin{cases}
    x = c_1 e^{-t} + c_2 e^{-6t} \\
    y = \frac{1}{2}c_1 e^{-t} - 2c_2 e^{-6t}
\end{cases}
\]
法二: 尝试解\\
此法将用到打量矩阵知识\\
将$(x,y)$看作一个列向量, 则:
\[
\begin{bmatrix}
    x \\
    y
\end{bmatrix}' 
= 
\begin{bmatrix}
    -2 &  2 \\
     2 & -5
\end{bmatrix}
\begin{bmatrix}
    x \\
    y
\end{bmatrix}
\]
解:
\[
\begin{bmatrix}
    x \\
    y
\end{bmatrix}
=
c_1
\begin{bmatrix}
    1 \\
    1/2
\end{bmatrix}
e^{-t}
+
c_2
\begin{bmatrix}
    1 \\
    -2
\end{bmatrix}
e^{-6t}
\]
现利用"尝试解":
\[
\begin{bmatrix}
    x \\
    y
\end{bmatrix}
=
\begin{bmatrix}
    a_1 \\
    a_2
\end{bmatrix}
e^{\lambda t}
\]
代入方程组:
\[
\begin{bmatrix}
    x \\
    y
\end{bmatrix}'
=
\lambda 
\begin{bmatrix}
    a_1 \\
    a_2
\end{bmatrix}
e^{\lambda t}
=
\begin{bmatrix}
    -2 &  2 \\
     2 & -5
\end{bmatrix}
\begin{bmatrix}
    a_1 \\
    a_2
\end{bmatrix}
e^{\lambda t}
\]
\[
\therefore  \hspace*{2em}
\lambda 
\begin{bmatrix}
    a_1 \\
    a_2
\end{bmatrix}
=
\begin{bmatrix}
    -2 &  2 \\
     2 & -5
\end{bmatrix}
\begin{bmatrix}
    a_1 \\
    a_2
\end{bmatrix}
\]
现即将解微分方程问题简化为解代数方程问题:
\[
\begin{cases}
    \lambda a_1 = -2 a_1 + 2 a_2\\
    \lambda a_2 = 2 a_1 - 5 a_2
\end{cases}
\]
将$\lambda$视为参数, 改写:
\[
\begin{cases}
    \lambda (-2 - \lambda)a_1 + 2 a_2 = 0\\
    \lambda 2 a_1 +(-5 - \lambda) a_2 = 0
\end{cases}
\]
即需解其非平凡解, 当且仅当对应行列式为$0$:
\[
\begin{vmatrix}
    -2 - \lambda & 2 \\
    2 & -5 - \lambda
\end{vmatrix}
= 0
\]
即:
\[
\textcolor{red}{\lambda^2 + 7 \lambda + 6 = 0}
\]
此方程称为特征方程, 解:
\[ \lambda = -1 \hspace*{0.5em} or -6 \]
当$\lambda = -1$:
\[
\begin{cases}
    -a_1 + 2 a_2 = 0 \\
    2 a_1 -4 a_2 = 0
\end{cases}
\]
令$a_2 = 1,  a_1= 2$, 即 
$\begin{bmatrix}
    2 \\
    1
\end{bmatrix}$  解:
\[
\begin{bmatrix}
    2 \\
    1
\end{bmatrix}
e^{-t}
\]
当$\lambda = -6$:
\[
\begin{cases}
    4a_1 + 2 a_2 = 0 \\
    2 a_1 + a_2 = 0
\end{cases}
\]
令$a_1 = 1,  a_1= -2$, 即 
$\begin{bmatrix}
    1 \\
    -2
\end{bmatrix}$  解:
\[
\begin{bmatrix}
    1 \\
    -2
\end{bmatrix}
e^{-6t}
\]
综上所述, 通解:
\[
\tilde{c_1}
\begin{bmatrix}
    2 \\
    1
\end{bmatrix}
e^{-t}
+
\tilde{c_2}
\begin{bmatrix}
    1 \\
    -2
\end{bmatrix}
e^{-6t}
\]
与法一解形式无区别.

\subsubsection{\large 矩阵形式一般化}
现将上述解法过程进行一般化:
\[
\begin{bmatrix}
    x \\
    y
\end{bmatrix}'
=
\begin{bmatrix}
    a & b \\
    c & d
\end{bmatrix}
\begin{bmatrix}
    x \\
    y
\end{bmatrix}
\]
尝试解:
\[
\begin{bmatrix}
    a_1 \\
    a_2
\end{bmatrix}
e^{\lambda t}
\]
代入:
\[
\lambda
\begin{bmatrix}
    a_1 \\
    a_2
\end{bmatrix}
=
\begin{bmatrix}
    a & b \\
    c & d
\end{bmatrix}
\begin{bmatrix}
    a_1 \\
    a_2
\end{bmatrix}
\]
有:
\[
\begin{bmatrix}
    a-\lambda & b \\
    c & d-\lambda
\end{bmatrix}
\begin{bmatrix}
    a_1 \\
    a_2
\end{bmatrix}
= 0
\]
求非平凡解, 当且仅当对应系数行列式为$0$:
\[
\begin{vmatrix}
    a-\lambda & b \\
    c & d-\lambda
\end{vmatrix}
= 0
\]
即:
\[
\textcolor{red}{\lambda^2 - (a+d)\lambda + (ad-bc) = 0}
\]
观察此特征方程\\
记
$A=
\begin{bmatrix}
    a & b \\
    c & d
\end{bmatrix}$, 有:
\[
tr(A)=a+d  \hspace*{2em} det(A)=ad-bc
\]
得特征方程:
\[
\lambda^2 - tr(A)\lambda + det(A) = 0
\]
假设其解为$\lambda_1, \lambda_2$(不同实数), 称为特征值.\\
对$\forall \lambda_i$ 找对应特征向量$\vec{\alpha _i}
=
\begin{bmatrix}
    a_{1i} \\
    a_{2i}
\end{bmatrix}$:
\[
\begin{bmatrix}
    a-\lambda_i & b \\
    c & d-\lambda_i
\end{bmatrix}
\begin{bmatrix}
    a_{1i} \\
    a_{2i}
\end{bmatrix}
= 0
\]
得通解:
\[
\begin{bmatrix}
    x \\
    y
\end{bmatrix}
=
c_1
\begin{bmatrix}
    a_{11} \\
    a_{21}
\end{bmatrix}
e^{\lambda_1 t}
+
c_2
\begin{bmatrix}
    a_{12} \\
    a_{22}
\end{bmatrix}
e^{\lambda_2 t}
\]

\subsubsection{\large 矩阵形式简化}
现继续对一般化过程简化, 可适用于$n \times n$:\\
记
$\vec{x} =
\begin{bmatrix}
    x \\
    y
\end{bmatrix}$ , 
$\vec{\alpha} =
\begin{bmatrix}
    a_1 \\
    a_2
\end{bmatrix}$ , 
$A =
\begin{bmatrix}
    a & b \\
    c & d
\end{bmatrix}$ \\
方程为:
\[
\vec{x}' = A \vec{x}
\]
尝试解:
\[
\vec{x}
=
\vec{\alpha}
e^{\lambda t}
\]
即:
\[ \lambda \vec{\alpha}= A \vec{\alpha} \]
\[ (A - \lambda I) \vec{\alpha} = 0 \]
特征方程:
\[
\begin{vmatrix}
    A - \lambda I
\end{vmatrix}
= 0
\]
根仍为特征值.

\subsubsection{\large 重根分析}
\[
\begin{vmatrix}
    A - \lambda I
\end{vmatrix}
= 0
\]
考虑到求解特征方程可能存在重根, 下举一例来分析如何处理:\\
eg: 
$A =
\begin{bmatrix}
    -2 & 1 & 1 \\
      1 & -2 & 1 \\
      1 & 1 & -2
\end{bmatrix}$\\
求解:
\[
\begin{vmatrix}
    A - \lambda I
\end{vmatrix}
=
\begin{vmatrix}
    -2-\lambda & 1 & 1 \\
      1 & -2-\lambda & 1 \\
      1 & 1 & -2-\lambda
\end{vmatrix}=0
\]
即得:
\[ \lambda (\lambda + 3)^2 = 0 \]
\[ \lambda = 0, \lambda = -3 (\text{重根})\]
当$\lambda = 0$时:
\[
\begin{cases}
    -2a_1 + a_2 + a_3 = 0 \\
    a_1 - 2a_2 + a_3 = 0 \\
    a_1 + a_2 - 2a_3 = 0
\end{cases}
\]
解:
\[
\vec{\alpha}_3  
=
\begin{bmatrix}
    1 \\
    1 \\
    1
\end{bmatrix}
\]
当$\lambda = -3$时:
\[
\begin{cases}
    a_1 + a_2 + a_3 = 0 \\
    a_1 + a_2 + a_3 = 0 \\
    a_1 + a_2 + a_3 = 0
\end{cases}
\]
令$a_1 = 1, a_2 = 0, a_3 = -1$, 即
\[
\vec{\alpha}_1 =
\begin{bmatrix}
    1 \\
    0 \\
    -1
\end{bmatrix}
e^{-3t}
\]
令$a_1 = 1, a_2 = -1, a_3 = 0$, 即
\[
\vec{\alpha}_2 =
\begin{bmatrix}
    1 \\
    -1 \\
    0
\end{bmatrix}
e^{-3t}
\]
通解:
\[
\vec{x}
=
c_1
\begin{bmatrix}
    1 \\
    0 \\
    -1
\end{bmatrix}
e^{-3t}
+
c_2
\begin{bmatrix}
    1 \\
    -1 \\
    0
\end{bmatrix}
e^{-3t}
+
c_3
\begin{bmatrix}
    1 \\
    1 \\
    1
\end{bmatrix}
\]
综上, 给出如下概念:\\
若 $\lambda$为重复特征值, 且能找到
足够的独立特征向量, 其构造所需数是独立解.
即称: 完备特征值.\\
若 不能找到足够的独立特征向量, 即称: 不完备特征值.\\[0.2em]
并给出如下定理:\\[0.5em]
主轴定理(谱定理):\\
若$A$为实对称 $n \times n$矩阵, 则它所有的特征值均是完备特征值.


\subsubsection{\large 复特征根}
到此, 考虑复特征值\\
遵循上节同一个程序, 仅改变定义域:
\[
\vec{\alpha } \cdot e^{(a + bi)t}
\]
在最后一步:\\
取实部和虚部, 得到两个实数解:
\[
\begin{cases}
    x' = x + 2y \\
    y' = -x - y
\end{cases}
, 
A =
\begin{bmatrix}
    1 & 2 \\
    -1 & -1
\end{bmatrix}
\]
得:
\[ \lambda^2 + 1 = 0 \]
\[ \lambda = \pm i \]
现令$a_1 = 1$, $a_2 = - \frac{1-i}{2}$
\[
\begin{bmatrix}
    1 \\
    \frac{i-1}{2}
\end{bmatrix}
e^{it} =
\left\{
\begin{bmatrix}
    1 \\
    -\frac{1}{2}
\end{bmatrix}
+
i
\begin{bmatrix}
    0 \\
    \frac{1}{2}
\end{bmatrix}
\right\} 
(\cos t + i \sin t)
\]
取$Re$ :
\[
\begin{bmatrix}
    1 \\
    -\frac{1}{2}
\end{bmatrix}
\cos t -
\begin{bmatrix}
    0  \\
    \frac{1}{2}
\end{bmatrix}
\sin t =
\begin{bmatrix}
    x \\
    y
\end{bmatrix}
\]
即:
\[
\begin{cases}
    x = \cos t \\
    y = -\frac{1}{2}(\sin t + \cos t)
\end{cases}
\]

\newpage
\subsection{\LARGE \hspace*{0.1em}2X2 常系数齐次线性方程组作图 }

\[
\begin{bmatrix}
    x \\
    y
\end{bmatrix}'
=
\begin{bmatrix}
    a & b \\
    c & d
\end{bmatrix}
\begin{bmatrix}
    x \\
    y
\end{bmatrix}
\]
eg:
\[
\begin{cases}
    x' = -x + by \\
    y' = cx - 3y
\end{cases} , x,y \text{表示偏离}
\]
求解:
\[
\begin{cases}
    x' = -x + 2y \\
    y' = - 3y
\end{cases} , 
A =
\begin{bmatrix}
    -1 & 2 \\
    0 & -3
\end{bmatrix}
\]
\[ \lambda^2 + \lambda + 3 = 0 \]
\[ \lambda = -3, \lambda = -1  \]
得:
\[
\begin{bmatrix}
    x \\
    y
\end{bmatrix}
=
c_1
\begin{bmatrix}
    1 \\
    -1
\end{bmatrix}
e^{-3t}
+
c_2
\begin{bmatrix}
    1 \\
    0
\end{bmatrix}
e^{-t}
\]
\[
\begin{cases}
    x = c_1 e^{-3t} + c_2 e^{-t} \\
    y = -c_1 e^{-3t} 
\end{cases}
\]
画草图:\\
先找基本点
\[
\begin{cases}
    c_1 = \pm 1 , &c_2 = 0 \\
    c_1 = 0, &c_2 = \pm 1
\end{cases}
\]
eg: 当$t \to \infty$时:
\[
c_1
\begin{bmatrix}
    1 \\
    -1
\end{bmatrix}
e^{-3t}
+
c_2
\begin{bmatrix}
    1 \\
    0
\end{bmatrix}
e^{-t}
\to 0
\]
作草图如下:\\
观察, 对于解汇聚于0(此例题), 称为汇聚节点(渐近稳定点)\\
若情况相反, 称为源节点(不稳定点).\\
\begin{figure}[htbp]
    \centering
\includegraphics[width=0.5\textwidth]{C:/Users/37/Desktop/Notes on Common Calculus Equations/Notes and pictures/cwf11.3.1.png}
    \caption{汇聚节点图}
    \label{fig:cwf11.3.1}
\end{figure}

再看下eg:
\[
\begin{cases}
    x' = -x + 3y \\
    y' = 5x - 3y
\end{cases} , 
A =
\begin{bmatrix}
    -1 & 3 \\
    5 & -3
\end{bmatrix}
\]
解得:
\[
\begin{bmatrix}
    x \\
    y
\end{bmatrix}
=
c_1
\begin{bmatrix}
    3 \\
    -5
\end{bmatrix}
e^{-6t}
+
c_2
\begin{bmatrix}
    1 \\
    1
\end{bmatrix}
e^{2t}
\]
\begin{figure}[htbp]
    \centering
\includegraphics[width=0.5\textwidth]{C:/Users/37/Desktop/Notes on Common Calculus Equations/Notes and pictures/cwf11.3.2.png}
    \caption{不稳定点图}
    \label{fig:cwf11.3.2}
\end{figure}

观察此解草图, 汇聚和源都不合适来形容, 仅称不稳定. \\
补eg:
\[
\begin{cases}
    x' = -x - y \\
    y' = 2x - 3y
\end{cases} , 
A =
\begin{bmatrix}
    -1 & -1 \\
    2 & -3
\end{bmatrix}
\]
\[ \lambda^2 + 4 \lambda + 5 = 0 \]
解得:
\[
\begin{bmatrix}
    x \\
    y
\end{bmatrix}
=
c_1
\left\{
\begin{bmatrix}
    a_1 \\
    a_2
\end{bmatrix}
\cos t +
\begin{bmatrix}
    b_1 \\
    b_2
\end{bmatrix}
\sin t
\right\} 
e^{-2t}
+
c_2
\ldots 
\]
显然, 此方程组图像有界.\\
有如下定理:\\
对下类方程:
\[ A x^2 + B y^2 + C xy + D x + E y + F = 0 \]
图像仅呈现 双曲线, 椭圆, 抛物线, 退化直线.\\
其中 只有椭圆有界.
\begin{figure}[htbp]
    \centering
\includegraphics[width=0.5\textwidth]{C:/Users/37/Desktop/Notes on Common Calculus Equations/Notes and pictures/cwf11.3.3.png}
    \caption{螺旋图}
    \label{fig:cwf11.3.3}
\end{figure}

其呈现螺旋汇聚现象(逆时针).
